\section{Botnet DDoS}
\label{sec:botnet}

The same ops portal targeted in
Section~\ref{sec:threat-dos} is the attack target. The distributed scenario is
simulated by cloning the attacker Incus container $N$ times, assigning each
clone a distinct IP address from the Incus bridge, and launching all bots
simultaneously. Each bot runs 20 concurrent threads for 90 seconds, easily
sustaining more than 50 requests per 10 seconds individually. The scenario is
demonstrated by \texttt{scripts/demo\_botnet.sh}, which accepts $N$ as an
argument. In testing, $N=3$ bots were used.

\subsection{Detection}

Detection extends the two-rule chain from Section~\ref{sec:threat-dos} with a
meta-rule that identifies the distributed pattern. Rule 100002 fires per bot,
as in the single-source case. Rule 100003 is a frequency rule that fires when
3 or more distinct source IPs each trigger rule 100002 within a 2-minute window.
The \texttt{<different\_field>client\_ip</different\_field>} primitive ensures
that the counter increments only when each triggering rule 100002 event came
from a new source IP. A single fast attacker cannot accidentally trigger it.

\begin{lstlisting}[style=xmlstyle, caption={Wazuh rule for botnet DDoS detection.}]
<rule id="100003" level="15" frequency="3"
      timeframe="120">
  <if_matched_sid>100002</if_matched_sid>
  <different_field>client_ip</different_field>
  <description>Botnet DDoS: 3+ distinct source IPs
    each exceeding flood threshold within
    2 minutes</description>
  <group>ddos,botnet,attack,</group>
  <options>alert_by_email</options>
</rule>
\end{lstlisting}

The 2-minute timeframe for rule 100003 is wider than the 10-second window of
rule 100002 to account for staggered bot startup times and jitter in network
delivery.

\subsection{Response}

Each bot's IP is blocked individually when rule 100002 fires for that source, via
the active response chain described in Section~\ref{sec:monitoring}. Rule 100003
does not trigger an additional block (the per-bot blocks from rule 100002
already handle individual sources). Rule 100003 provides the distributed-attack
signal: its level-15 alert notifies operators that a coordinated attack is in
progress, as distinct from isolated single-source incidents.

\subsection{Evaluation}

In testing with 3 bots, rule 100002 fired for each bot within the first
10-second window of their respective floods, and rule 100003 fired 12 times over
the 90-second run as new combinations of three distinct IPs were detected. The
Grafana \emph{Botnet Patterns Detected} stat panel (rule 100003 count) and
\emph{Distinct Bot IPs} stat (count of distinct IPs triggering rule 100002)
confirm the distributed pattern. The \emph{Bot Flood - Threshold Breaches per
Source IP} timeseries displays one line per bot IP, with simultaneous spikes
constituting the visual signature of a coordinated flood.

The principal limitation is that rule 100003 requires each individual bot to
exceed the 50-request-per-10-second threshold. A large botnet coordinating a
slow distributed flood (where each bot sends just below 50 requests per 10
seconds, but their aggregate rate is sufficient to degrade service) would not
trigger rule 100002 per bot, and consequently rule 100003 would not fire.
Detecting such a low-rate distributed pattern would require aggregate rate
monitoring across all source IPs, which is not implemented in the current rule
set.
